\documentclass[11pt]{article}

\usepackage[a4paper,margin=1in]{geometry}
\usepackage{amsmath,amssymb}
\usepackage{hyperref}
\usepackage{tabularx}
\usepackage{booktabs}
\usepackage{microtype}

\title{p-Adic Numbers, Valuations, and Quantum Circuits:\\
A Focused Survey for Multiplicity-Based Quantum Pipelines}

\author{Citizen Gardens \\ The Prime Materia Commons}
\date{May 5, 2026}

\begin{document}
\maketitle

\begin{abstract}
We survey a small but conceptually rich literature connecting p-adic numbers
and valuations to quantum theory, statistical systems, and computational
tooling. Our emphasis is on work that is structurally relevant to
multiplicity-based, valuation-driven feedback loops in noisy
intermediate-scale quantum (NISQ) pipelines such as M--QNN.
We highlight p-adic quantum mechanics, p-adic qubits and gates,
p-adic statistical systems, FPGA-based quantum emulation,
and classical p-adic arithmetic libraries, and we identify a gap:
no existing work realizes the kind of prime-exponent valuation dynamics
we engineer in M--QNN as part of an adaptive measurement and training loop.
\end{abstract}

\section{Introduction and Motivation}

Classical p-adic number theory encodes multiplicities of prime factors via
the p-adic valuation \(v_p\), which measures the exponent of a prime \(p\)
in the factorization of an integer or, more generally, a number field element.
Over the last several decades, p-adic structures have been imported into
mathematical physics and quantum theory, leading to p-adic Hilbert spaces,
p-adic wave functions, and p-adic statistical ensembles
(e.g.~Khrennikov's overview~\cite{khrennikov2017padic},
Aniello--Mancini--Parisi~\cite{aniello2023trace},
and subsequent work on p-adic quantum mechanics).

In parallel, noisy intermediate-scale quantum (NISQ) algorithms such as
M--QNN use standard complex-valued quantum circuits but rely heavily
on classical bookkeeping of measurement multiplicities and statistical
confidence bounds.
From a multiplicity-theoretic point of view, these discrete counts
are exponents of primes in a structured state space.  This suggests
a conceptual bridge between p-adic valuations and the engineered
feedback loops of NISQ-era quantum machine learning:
prime-exponent valuations as state variables, with adaptive update
rules and stopping inequalities defined directly on the valuation lattice.

This note collects the most relevant post-2015 work that touches this
bridge from the p-adic side, and contrasts it with the multiplicity-driven
design used in M--QNN.

\section{Core Works and Their Relation to Multiplicity-Based Pipelines}

Table~\ref{tab:core} summarizes six representative works and toolchains
that are most relevant to a p-adic, valuation-centric view of quantum
circuits and statistical systems.

\begin{table}[h]
\centering
\small
\renewcommand{\arraystretch}{1.2}
\begin{tabularx}{\textwidth}{@{}p{2.3cm}p{2.2cm}p{2.2cm}p{1.8cm}X@{}}
\toprule
\textbf{Short citation} &
\textbf{p-adic use} &
\textbf{Quantum model} &
\textbf{Hardware link} &
\textbf{Main method, key result, and relevance to multiplicity view} \\
\midrule

Khrennikov et al.~2017
(\emph{p-Adic Mathematical Physics: The First 30 Years})
&
p-adic Hilbert spaces; wave functions and amplitudes over \(\mathbb{Q}_p\)
&
Abstract p-adic quantum mechanics (states, observables, dynamics)
&
None
&
Reconstructs quantum-mechanical structures (states, observables,
probability amplitudes) over non-Archimedean fields, showing that a
consistent p-adic quantum formalism exists.
Conceptually supports treating valuations and p-adic fields as legitimate
state spaces, aligning with prime-exponent lattices as structured
multiplicity spaces. \\

Aniello--Mancini--Parisi 2023
(\emph{Trace class operators and states in p-adic QM})
&
States and trace-class operators over quadratic extensions of \(\mathbb{Q}_p\);
p-adic probability weights
&
Rigorous operator-theoretic p-adic quantum mechanics
&
None
&
Develops functional analysis for p-adic Banach spaces, defines trace,
density operators, and expectation values in p-adic quantum mechanics.
Provides formal tools (trace, density matrices) that could describe
probability distributions over multiplicity-encoded or prime-exponent
states in a valuation-centric model. \\

Svampa et al.~2026
(\emph{Composing p-adic qubits: from representations of SO(3)\(_p\) \dots})
&
p-adic qubits as irreducible representations of SO(3)\(_p\); p-adic
controlled gates
&
p-adic qubit and gate model, including entanglement
&
Implicit (gate-level architecture; not mapped to physical qubits)
&
Uses representation theory of compact p-adic groups to define p-adic
qubits and multi-qubit spaces, and constructs universal p-adic gate sets
(explicitly analyzed for \(p=3\)).
Offers a natural formalism for a ``native'' p-adic circuit picture of
multiplicity-state evolution, in contrast to classical valuation
bookkeeping in NISQ pipelines. \\

Dragovich et al.~2021
(\emph{Statistical system based on p-adic numbers})
&
Configuration space and phase variables in \(\mathbb{Q}_p\);
p-adic partition functions
&
None (classical statistical mechanics)
&
None
&
Builds statistical ensembles with p-adic configuration spaces and
defines partition functions as p-adic sums, then maps to real-valued
thermodynamic quantities (free energy, entropy, specific heat).
Demonstrates how valuation-structured state spaces can produce
real-valued observables, analogous to multiplicity-based stopping
criteria (e.g.\ Hoeffding bounds) in M--QNN. \\

ManishPatla 2023
(\emph{QuantumComputation\_FPGAs} repository)
&
No explicit p-adics; fixed-point complex arithmetic pipelines
&
Gate-model quantum circuit emulation (state-vector)
&
FPGA (VHDL-based emulation of quantum gates)
&
Implements gate-level emulation of quantum circuits on FPGAs using
fixed-point arithmetic for amplitudes and models of quantum noise.
Provides a hardware blueprint for building dedicated controllers that
could update prime-exponent / valuation registers in lockstep with
quantum measurements in adaptive, multiplicity-driven NISQ pipelines. \\

SageMath p-adic valuation tools
(e.g.\ \texttt{padic\_valuation.py})
&
Full p-adic arithmetic: valuations on integers and number fields,
completions \(\mathbb{Q}_p\)
&
None (classical CAS and number-theory tooling)
&
Software library (Python/C)
&
Supplies robust algorithms for exact p-adic arithmetic, valuations,
and expansions in computer algebra.
Essentially the classical analogue of the valuation bookkeeping layer
in M--QNN-like feedback loops, showing that treating valuations as
first-class computational objects is standard and well-supported. \\

\bottomrule
\end{tabularx}
\caption{Core works linking p-adic numbers, valuations, and quantum or
quantum-adjacent computational models, with emphasis on relevance to
multiplicity-based NISQ pipelines such as M--QNN.}
\label{tab:core}
\end{table}

\section{Discussion: From p-Adic Models to Multiplicity-Based Feedback}

% Here you can connect the table to your M--QNN work:
% - valuations as exponents in a prime-indexed multiplicity lattice
% - early-exit rules as inequalities on valuations
% - gap between p-adic Hilbert-space models and NISQ-era hybrid pipelines.

\section{Conclusion and Open Directions}

% Summarize gaps and propose future work:
% - native p-adic circuit models vs. classical valuation control loops
% - hardware support (FPGA/ASIC) for valuation updates
% - integration with variational training and adaptive measurement.

%=========================================================
% Mathematical Appendix: Valuation-Based Early-Exit and
% Operator Bounds for M--QNN-Style Feedback Loops
%=========================================================

\appendix
\section{Mathematical Appendix}
\label{app:math}

This appendix collects explicit proofs and norm bounds underlying the
valuation-based early-exit rule and related operators in the M--QNN
pipeline.  We work in a simplified setting (two and then many
candidates), but all arguments extend to the full candidate set used in
practice.

\subsection{Setup and Notation}

Fix a query \(x\) and a finite reference set of candidates
\(j \in \mathcal{C}(x)\).
For each candidate \(j\) we estimate the fidelity
\(F_j := F(x,x_j)\in[0,1]\) using the swap test.
One swap-test shot for candidate \(j\) produces an ancilla outcome
\(Z_{j,s}\in\{0,1\}\) with
\begin{equation}
  \mathbb{P}(Z_{j,s}=0)
  \;=\;
  p_{j,0}
  \;:=\;
  \frac{1+F_j}{2},
  \qquad
  \mathbb{P}(Z_{j,s}=1) = 1-p_{j,0}.
\end{equation}
We assume conditional independence of shots:
\(\{Z_{j,s}\}_{s\ge 1}\) are i.i.d.\ for each fixed \(j\),
and independent across different \(j\).

After \(S_j\) shots for candidate \(j\), define the multiplicities
\begin{equation}
  N_{j,0} = \sum_{s=1}^{S_j} \mathbf{1}[Z_{j,s}=0],
  \qquad
  N_{j,1} = S_j - N_{j,0},
\end{equation}
the empirical probability
\(\hat{p}_{j,0} := N_{j,0}/S_j\),
and the fidelity estimator
\begin{equation}
  \hat{F}_j
  :=
  2\hat{p}_{j,0} - 1
  \;=\;
  2\frac{N_{j,0}}{S_j} - 1.
\end{equation}
For a confidence parameter \(\delta \in (0,1)\), define the
Hoeffding radius
\begin{equation}
  r_j
  :=
  \sqrt{\frac{\ln(2/\delta)}{2 S_j}}.
\end{equation}
The classical (scalar) early-exit rule declares candidate \(j^\star\) as
winner when
\begin{equation}
  \hat{F}_{j^\star} - r_{j^\star}
  >
  \max_{j\neq j^\star} \bigl(\hat{F}_j + r_j\bigr).
  \label{eq:early-exit-app}
\end{equation}

\subsection{Explicit multiplicity inequality (two-candidate case)}
\label{app:two-candidate}

We first specialize \eqref{eq:early-exit-app} to two candidates
\(j\in\{1,2\}\) and derive the early-exit condition as an explicit
inequality in the integer multiplicities
\((N_{1,0},N_{2,0},S_1,S_2)\).

\begin{lemma}[Explicit multiplicity inequality, two candidates]
\label{lem:mult-ineq-two}
In the two-candidate setting, the condition
\begin{equation}
  \hat{F}_1 - r_1 > \hat{F}_2 + r_2
\end{equation}
is equivalent to
\begin{equation}
  2\left(\frac{N_{1,0}}{S_1} - \frac{N_{2,0}}{S_2}\right)
  >
  \sqrt{\frac{\ln(2/\delta)}{2}}
  \left(
    \frac{1}{\sqrt{S_1}}
    +
    \frac{1}{\sqrt{S_2}}
  \right),
  \label{eq:mult-ineq-two}
\end{equation}
which is an explicit inequality in the integer multiplicities.
\end{lemma}

\begin{proof}
Write the early-exit condition for candidate~1 as
\begin{equation}
  \left(2\frac{N_{1,0}}{S_1} - 1\right)
  - \sqrt{\frac{\ln(2/\delta)}{2 S_1}}
  >
  \left(2\frac{N_{2,0}}{S_2} - 1\right)
  + \sqrt{\frac{\ln(2/\delta)}{2 S_2}}.
\end{equation}
Adding \(1\) to both sides cancels the \(-1\) terms and yields
\begin{equation}
  2\frac{N_{1,0}}{S_1}
  - \sqrt{\frac{\ln(2/\delta)}{2 S_1}}
  >
  2\frac{N_{2,0}}{S_2}
  + \sqrt{\frac{\ln(2/\delta)}{2 S_2}}.
\end{equation}
Rearranging,
\begin{equation}
  2\left(\frac{N_{1,0}}{S_1} - \frac{N_{2,0}}{S_2}\right)
  >
  \sqrt{\frac{\ln(2/\delta)}{2 S_1}}
  + \sqrt{\frac{\ln(2/\delta)}{2 S_2}}.
\end{equation}
Factoring out \(\sqrt{\ln(2/\delta)/2}\) on the right-hand side,
\begin{equation}
  2\left(\frac{N_{1,0}}{S_1} - \frac{N_{2,0}}{S_2}\right)
  >
  \sqrt{\frac{\ln(2/\delta)}{2}}
  \left(
    \frac{1}{\sqrt{S_1}}
    +
    \frac{1}{\sqrt{S_2}}
  \right),
\end{equation}
which is precisely \eqref{eq:mult-ineq-two}.
\end{proof}

For some purposes it is convenient to express the left-hand side as a
\emph{cross-multiplicity} term:
\begin{equation}
  \Delta
  :=
  N_{1,0} S_2 - N_{2,0} S_1.
\end{equation}
Multiplying both sides of \eqref{eq:mult-ineq-two} by \(S_1 S_2\) gives
\begin{equation}
  2 \Delta
  >
  \sqrt{\frac{\ln(2/\delta)}{2}}
  \left(
    S_2 \sqrt{S_2} + S_1 \sqrt{S_1}
  \right),
\end{equation}
which isolates an integer quantity \(\Delta\) on the left and a
scale-dependent bound on the right.

\subsection{High-probability correctness of early exit}
\label{app:hoeffding}

We now prove that the early-exit rule \eqref{eq:early-exit-app}
identifies the true best candidate with probability at least \(1-\delta\)
under the independence assumptions stated above.

\begin{lemma}[Hoeffding interval for fidelity estimator]
\label{lem:hoeffding}
For each candidate \(j\), we have
\begin{equation}
  \mathbb{P}\Bigl(
    \bigl|\hat{F}_j - F_j\bigr|
    \le r_j
  \Bigr)
  \;\ge\;
  1-\delta,
\end{equation}
where \(r_j = \sqrt{\ln(2/\delta)/(2 S_j)}\).
\end{lemma}

\begin{proof}
Since \(Z_{j,s}\in\{0,1\}\) are i.i.d.\ with mean \(p_{j,0}\),
Hoeffding's inequality for bounded variables on \([0,1]\) gives
\begin{equation}
  \mathbb{P}\Bigl(\bigl|\hat{p}_{j,0} - p_{j,0}\bigr| \ge t\Bigr)
  \le 2 \exp(-2 S_j t^2)
  \quad
  \forall t>0.
\end{equation}
Set \(t = \sqrt{\ln(2/\delta)/(2 S_j)}\); then
\begin{equation}
  2 \exp(-2 S_j t^2)
  =
  2 \exp\left(
    -2 S_j \cdot \frac{\ln(2/\delta)}{2 S_j}
  \right)
  =
  2 \exp\bigl(-\ln(2/\delta)\bigr)
  =
  \delta.
\end{equation}
Thus
\begin{equation}
  \mathbb{P}\Bigl(
    \bigl|\hat{p}_{j,0} - p_{j,0}\bigr|
    \le t
  \Bigr)
  \ge 1-\delta.
\end{equation}
Because \(\hat{F}_j = 2\hat{p}_{j,0}-1\) and \(F_j = 2 p_{j,0}-1\),
\begin{equation}
  \bigl|\hat{F}_j - F_j\bigr|
  =
  2 \bigl|\hat{p}_{j,0} - p_{j,0}\bigr|.
\end{equation}
Therefore
\begin{equation}
  \mathbb{P}\Bigl(
    \bigl|\hat{F}_j - F_j\bigr|
    \le 2 t
  \Bigr)
  \ge 1-\delta.
\end{equation}
By definition, \(r_j = 2 t\), hence the desired inequality.
\end{proof}

We now show that the early-exit rule correctly identifies the best
candidate with high probability.

\begin{theorem}[Correctness of early exit]
\label{thm:correct-early-exit}
Let \(j^\star\) be the true best candidate, i.e.,
\(F_{j^\star} = \max_j F_j\).
Assume that for each candidate \(j\), the fidelity estimator
\(\hat{F}_j\) is built from i.i.d.\ swap-test shots as above.
If the early-exit condition \eqref{eq:early-exit-app} holds for some
index \(\hat{j}\), then under the event
\(\mathcal{E} = \bigcap_j \{\lvert \hat{F}_j - F_j\rvert \le r_j\}\)
we must have \(\hat{j} = j^\star\).
Moreover,
\begin{equation}
  \mathbb{P}(\hat{j} \neq j^\star)
  \le
  \mathbb{P}(\mathcal{E}^c),
\end{equation}
and if the \(\hat{F}_j\) are conditionally independent across \(j\), then
\begin{equation}
  \mathbb{P}(\hat{j} \neq j^\star)
  \le
  \sum_j \mathbb{P}\bigl(\lvert \hat{F}_j - F_j\rvert > r_j\bigr)
  \le
  |\mathcal{C}(x)| \,\delta.
\end{equation}
\end{theorem}

\begin{proof}
On the event \(\mathcal{E}\) we have, for every \(j\),
\begin{equation}
  F_j
  \ge \hat{F}_j - r_j,
  \qquad
  F_j
  \le \hat{F}_j + r_j.
\end{equation}
Suppose the early-exit condition holds with index \(\hat{j}\):
\begin{equation}
  \hat{F}_{\hat{j}} - r_{\hat{j}}
  >
  \max_{j\neq \hat{j}} (\hat{F}_j + r_j).
  \label{eq:ee-proof-cond}
\end{equation}
Fix any \(k\neq \hat{j}\).  Using the event \(\mathcal{E}\),
\begin{equation}
  F_{\hat{j}}
  \ge \hat{F}_{\hat{j}} - r_{\hat{j}},
  \qquad
  F_k
  \le \hat{F}_k + r_k.
\end{equation}
Combining these inequalities with \eqref{eq:ee-proof-cond},
\begin{equation}
  F_{\hat{j}}
  \ge \hat{F}_{\hat{j}} - r_{\hat{j}}
  >
  \hat{F}_k + r_k
  \ge F_k.
\end{equation}
Thus \(F_{\hat{j}} > F_k\) for every \(k\neq \hat{j}\), i.e.,
\(\hat{j}\) attains the unique maximum of the true fidelities, so
\(\hat{j}=j^\star\).

For the error probability, note that under \(\mathcal{E}\) we can never
output the wrong index, hence
\(\{\hat{j} \neq j^\star\} \subseteq \mathcal{E}^c\),
so
\begin{equation}
  \mathbb{P}(\hat{j} \neq j^\star)
  \le
  \mathbb{P}(\mathcal{E}^c).
\end{equation}
If the estimators are independent across \(j\), then by the union bound
and Lemma~\ref{lem:hoeffding},
\begin{equation}
  \mathbb{P}(\mathcal{E}^c)
  \le
  \sum_{j\in \mathcal{C}(x)}
  \mathbb{P}\bigl(\lvert\hat{F}_j - F_j\rvert > r_j\bigr)
  \le
  |\mathcal{C}(x)| \,\delta,
\end{equation}
which completes the proof.
\end{proof}

\begin{remark}[Scaled \(\delta\) choice]
In practice one often chooses a per-candidate confidence level
\(\delta' := \delta / |\mathcal{C}(x)|\) and defines
\(r_j = \sqrt{\ln(2/\delta')/(2 S_j)}\).
By Lemma~\ref{lem:hoeffding} and the union bound we then obtain
\(\mathbb{P}(\mathcal{E}^c)\le \delta\), so the probability of selecting
an incorrect winner is at most \(\delta\).
\end{remark}

\subsection{Operator-norm bounds for fidelity estimators}
\label{app:operator-norm}

We now recast the scalar Hoeffding bound in terms of operator norms of
the swap-test measurement operator, which is useful when viewing the
entire estimation process as a linear map on density operators.

\subsubsection{Swap-test measurement operator}

For a fixed candidate pair \((x,x_j)\), the swap test can be described
by a POVM \(\{M_0, M_1\}\) on the joint query--reference--ancilla
Hilbert space \(\mathcal{H}_Q \otimes \mathcal{H}_R \otimes \mathcal{H}_A\)
such that
\begin{equation}
  \mathrm{Pr}[Z=0]
  =
  \mathrm{Tr}(M_0 \rho),
  \qquad
  \mathrm{Pr}[Z=1]
  =
  \mathrm{Tr}(M_1 \rho),
\end{equation}
with \(\rho = \rho_Q \otimes \rho_R \otimes \rho_A\) the joint state and
\(M_0 + M_1 = I\).
For pure states and appropriate initialization, one obtains
\(\mathrm{Pr}[Z=0] = (1+F_j)/2\).

Define the centered observable
\begin{equation}
  O
  :=
  2 M_0 - I,
\end{equation}
so that for a single shot \(Z\),
\begin{equation}
  \mathbb{E}[2\mathbf{1}[Z=0]-1]
  =
  \mathrm{Tr}(O \rho)
  =
  F_j.
\end{equation}
The empirical estimator after \(S\) shots is
\begin{equation}
  \hat{F}
  =
  \frac{1}{S} \sum_{s=1}^S X_s,
  \qquad
  X_s := 2\mathbf{1}[Z_s=0]-1,
\end{equation}
with \(\mathbb{E}[X_s] = \mathrm{Tr}(O\rho)\).

\begin{lemma}[Operator norm of \(O\)]
\label{lem:op-norm}
The operator \(O = 2M_0 - I\) has operator norm
\(\|O\|_{\mathrm{op}} \le 1\).
\end{lemma}

\begin{proof}
The POVM elements satisfy \(0 \preceq M_0 \preceq I\), so all
eigenvalues \(\lambda_k(M_0)\) lie in \([0,1]\).
The eigenvalues of \(O\) are \(2\lambda_k(M_0) - 1\), which lie in
\([-1,1]\).
Hence \(\|O\|_{\mathrm{op}} = \max_k |2\lambda_k(M_0)-1| \le 1\).
\end{proof}

\subsubsection{Variance and concentration in operator form}

We can view each \(X_s\) as a random variable with values in \([-1,1]\)
and mean \(\mathbb{E}[X_s] = \mathrm{Tr}(O\rho)\).
Then\footnote{This is the same variance bound used implicitly in the
main text; we spell it out here for completeness.}
\begin{equation}
  \mathrm{Var}(X_s)
  \le 1,
  \qquad
  \mathrm{Var}(\hat{F})
  =
  \mathrm{Var}\Bigl(\frac{1}{S}\sum_{s=1}^S X_s\Bigr)
  =
  \frac{\mathrm{Var}(X_s)}{S}
  \le \frac{1}{S}.
\end{equation}
Hoeffding's inequality for bounded variables on \([-1,1]\) yields
\begin{equation}
  \mathbb{P}\Bigl(
    |\hat{F} - F_j|
    \ge t
  \Bigr)
  \le 2\exp\Bigl(-\frac{S t^2}{2}\Bigr),
\end{equation}
which matches the scalar bound of Lemma~\ref{lem:hoeffding} after
reparametrization.

\subsection{Valuation-based encoding and inequalities}
\label{app:valuation}

For completeness, we record the valuation-based encoding of the
multiplicities and the equivalence of the multiplicity inequality with
a valuation inequality.

Fix distinct primes \(p_j, q_j\) for each candidate \(j\).
Define the global integer
\begin{equation}
  M
  :=
  \prod_{j\in\mathcal{C}(x)}
  p_j^{S_j} \, q_j^{N_{j,0}}.
\end{equation}
The \(p_j\)-adic and \(q_j\)-adic valuations are
\begin{equation}
  v_{p_j}(M) = S_j,
  \qquad
  v_{q_j}(M) = N_{j,0}.
\end{equation}

In the two-candidate case, inequality
\eqref{eq:mult-ineq-two} can be rewritten as
\begin{equation}
  2\left(
    \frac{v_{q_1}(M)}{v_{p_1}(M)}
    -
    \frac{v_{q_2}(M)}{v_{p_2}(M)}
  \right)
  >
  \sqrt{\frac{\ln(2/\delta)}{2}}
  \left(
    \frac{1}{\sqrt{v_{p_1}(M)}}
    +
    \frac{1}{\sqrt{v_{p_2}(M)}}
  \right),
\end{equation}
which shows that the early-exit decision can be expressed as a
predicate on the valuations of \(M\).
This viewpoint is useful in multiplicity-based models where prime
exponents are treated as state coordinates.

\medskip

The arguments in this appendix provide a mathematically explicit
foundation for the valuation-based early-exit rule and the associated
operator bounds used in M--QNN and related NISQ-practical pipelines.

\nocite{*}
\bibliographystyle{unsrtnat}
\bibliography{references}
\end{document}
